\problem{CA01}{数据流中的加权抽样}{}
数据项按顺序到达，第 $i$ 项带正权重 $w_i$。事先不知道总项数，只允许保存一个候选项和常数个数值。设计算法，使读完任意 $t\ge1$ 项后，候选为第 $i$ 项的概率恰为 $w_i/\sum_{j=1}^t w_j$。算法每步可使用一个独立的连续均匀随机数。证明结论，并说明为什么“新项比旧候选重就替换”不正确。
\solution 维护累计权重 $W$。第一项到达时令候选为它、$W=w_1$。第 $t$ 项到达时先令 $W\leftarrow W+w_t$，再以概率 $w_t/W$ 替换候选，否则保留。一次更新只需常数次运算，保存候选标识及累计权重即可。

用归纳证明不变量。设前 $t-1$ 项累计权重为 $W_{t-1}$。旧项 $i<t$ 成为新候选，需要它原来被选中且这次未被替换；当前随机数与过去独立，故
\[
\Prob(\text{候选为 }i)
=\frac{w_i}{W_{t-1}}\left(1-\frac{w_t}{W_t}\right)
=\frac{w_i}{W_t}.
\]
新项成为候选的概率则直接为 $w_t/W_t$。归纳起点 $t=1$ 显然成立，因此对每个确定的 $t$ 都满足要求；结束长度由外部给定且不依赖抽样结果时，结论也成立。

例如权重依次为 $1,2,3$，第二步选第二项的概率为 $2/3$，第三步替换概率为 $1/2$，最终三项概率为 $1/6,1/3,1/2$。按候选权重大小作确定性替换，在严格递增的权重流中总会保留最后一项，无法产生所需分布。

算法对固定权重序列成立；若权重随机，但与抽样所用随机数独立，可以先条件于整条权重序列应用同一证明。若后续权重由已公开的候选结果自适应决定，不能直接沿用这个条件化结论。空间界按保存的数值个数计算，精确累计权重所需的位数可能增长。
\insight 保存的是满足概率不变量的代表项。新项应与全部历史权重竞争，而不是只与当前候选的权重比较。

\problem{CA02}{两个有序数组的中位数}{}
给定两个非递减实数数组 $A,B$，长度分别为 $m,n$，且 $m+n>0$。允许随机访问，不允许先合并全部元素。求合并样本的中位数：总数为偶数时取中间两数的平均值。设计 $O(\log(\min\{m,n\}+2))$ 时间、$O(1)$ 额外空间的算法，处理空数组、重复值和极不相同的数组长度。
\solution 必要时交换两数组，使 $m\le n$。若 $m=0$，直接读取 $B$ 的中央一项或两项。以下设 $m>0$，令
\[
h=\left\lfloor\frac{m+n+1}{2}\right\rfloor.
\]
从 $A$ 的前部取 $i$ 项、从 $B$ 的前部取 $j=h-i$ 项作为左半部。合法搜索区间为
\[
\max(0,h-n)\le i\le\min(m,h).
\]
定义分割两侧的边界元素
\[
\begin{array}{ll}
L_A=A_{i-1}\ (i>0),& R_A=A_i\ (i<m),\\
L_B=B_{j-1}\ (j>0),& R_B=B_j\ (j<n).
\end{array}
\]
缺失的左边界视为 $-\infty$，缺失的右边界视为 $+\infty$；这些哨兵仅用于比较。

正确分割的充要条件是 $L_A\le R_B$ 且 $L_B\le R_A$。因为数组各自有序，这两个交叉比较就保证整个左半部不大于整个右半部。若 $L_A>R_B$，说明从 $A$ 取多了，应减小 $i$；随着 $i$ 增大，$L_A$ 不减而 $R_B$ 不增，故所有更大的 $i$ 都不可能正确。若 $L_B>R_A$，同理应增大 $i$。据此在合法整数区间二分。

正确分割必存在：将全部元素概念上排序并对相等值固定任一稳定顺序，前 $h$ 项在每个原数组内都是前缀。二分的每次排除都保留至少一个正确分割，因此必终止。找到后，若 $m+n$ 为奇数，答案为 $\max(L_A,L_B)$；若为偶数，答案为
\[
\frac{\max(L_A,L_B)+\min(R_A,R_B)}2.
\]
总数为偶数且非零时两半均非空，因此最终平均的都是有限真实元素。重复值允许多个正确分割，但它们给出的中位数相同。搜索区间至多含 $m+1$ 个整数，得到所需时间和空间界。
\insight 二分搜索的对象是左半部应从较短数组取多少项。证明方向单调后，不必实际合并任何元素。

\problem{CA03}{逆序对与最少相邻交换}{}
实数数组 $a_1,\ldots,a_n$ 可有重复值。逆序对定义为 $i<j$ 且 $a_i>a_j$。设计 $O(n\log n)$ 时间算法计数，并证明将数组变成非递减顺序所需的最少相邻交换次数恰等于逆序对数。计算 $(3,1,2,1)$ 的结果，说明相等元素的处理。
\solution 用归并排序同时计算逆序对。先分别递归处理左右两半，累计内部逆序对；再将已排序的两半合并。若当前左值不大于右值，就输出左值。否则输出右值，且左半部所有尚未输出的值都严格大于该右值，故一次增加“左边剩余元素数”个跨半区逆序对。每对跨半区元素只在其右元素先被输出时计数一次，因此不遗漏、不重复。

每层合并总工作量为 $O(n)$，层数为 $O(\log n)$，额外缓冲区为 $O(n)$。空数组或单元素数组的计数为零。相等时优先输出左值，避免把相等元素误计为逆序；计数变量须能容纳 $n(n-1)/2$。

再证明交换数的最优性。交换一对相邻元素时，它们与其他元素形成的逆序对总数不变，因为其他元素位于这两项的同一侧；只有交换的这对关系可能改变。因此一次交换使总逆序数至多下降一。排序后的逆序数为零，所以任何方案至少需要初始逆序数次交换。

若数组尚未有序，必有相邻的严格下降对，否则所有相邻关系非递减就已保证整体有序。交换这样的下降对，逆序数恰减少一。重复操作，恰在初始逆序数步后到达有序数组，达到下界。例中逆序对为
\[
(1,2),(1,3),(1,4),(3,4),
\]
总数为 $4$，故最少相邻交换次数也是 $4$。
\insight 分治使每一对元素有唯一的计数时机；逆序数又是证明交换操作最优性的势函数。
